% !TEX TS-program = lualatexmk
%% %%%%%%%%%%%%%%%%%%%%%%%%%%%%%%%%%%%%%%%%%%%%%%%%%%%%%%%%%%%%%%%%%%%%
%% Erste-Schritte.tex – ein LaTeX-Dokument für absolute Anfänger
%% Stand: 2026-09-24, U. Groh
%%
%% So benutzt man diese Datei:
%%   1. Datei in Overleaf hochladen oder in TeXworks/TeXShop öffnen.
%%   2. Übersetzen („Compile“ bzw. „Setzen“) – am besten mit LuaLaTeX;
%%      pdfLaTeX geht auch.
%%   3. Für das Literaturverzeichnis muss im Editor „biber“ eingestellt
%%      sein (nicht „bibtex“). In Overleaf passiert das automatisch.
%%   4. Beim ersten Mal zwei- bis dreimal übersetzen, damit Inhalts-
%%      verzeichnis, Querverweise und Literaturangaben stimmen.
%%
%% Alles, was hinter einem Prozentzeichen % steht, ist ein Kommentar 
%% und erscheint NICHT im PDF. Diese Kommentare erklären die Datei.
%% %%%%%%%%%%%%%%%%%%%%%%%%%%%%%%%%%%%%%%%%%%%%%%%%%%%%%%%%%%%%%%%%%%%%

%% ====================================================================
%% TEIL 0: DIE LITERATURDATENBANK
%% Die folgenden Zeilen schreibt LaTeX beim Übersetzen in die Datei
%% Erste-Schritte.bib. Später nimmt man besser eine eigene .bib-Datei
%% und trägt ihren Namen unten bei \addbibresource ein.
%% Einträge findet man z. B. bei https://zbmath.org (Export „BibTeX“).
%% ====================================================================
\begin{filecontents*}[overwrite]{\jobname.bib}
@book{heuser,
  author    = {Heuser, Harro},
  title     = {Lehrbuch der Analysis. Teil 1},
  publisher = {Vieweg+Teubner},
  location  = {Wiesbaden},
  edition   = {17},
  year      = {2009}
}
@manual{lshort,
  author = {Daniel, Marco and Gundlach, Patrick and Schmidt, Walter},
  title  = {\LaTeXe-Kurzbeschreibung},
  url    = {https://ctan.org/pkg/lshort-german}
}
\end{filecontents*}

%% ====================================================================
%% TEIL 1: DIE PRÄAMBEL
%% Alles zwischen \documentclass und \begin{document} ist die Präambel.
%% Hier wird festgelegt, WIE das Dokument aussieht – nicht WAS drinsteht.
%% ====================================================================

%% -- Dokumentklasse: scrartcl (KOMA-Script) für Artikel, Hausarbeiten,
%%    Seminarausarbeitungen. Optionen stehen in eckigen Klammern.
\documentclass[%
	,ngerman				% Sprache Deutsch (neue Rechtschreibung)
	,fontsize	= 11pt		% Schriftgröße: 10pt, 11pt oder 12pt
	,parskip	= half		% Absätze durch eine halbe Leerzeile trennen
	]{scrartcl}

%% -- Schrift und Sprache
\usepackage{iftex}				% erkennt, ob pdfLaTeX oder LuaLaTeX läuft
\ifPDFTeX						% nur für pdfLaTeX nötig:
	\usepackage[T1]{fontenc}	%   richtige Umlaute und Silbentrennung
	\usepackage{lmodern}		%   Schrift Latin Modern
\fi
\usepackage{babel}				% deutsche Silbentrennung, „Inhaltsverzeichnis“ usw.
\usepackage{csquotes}			% richtige Anführungszeichen mit \enquote{...}

%% -- Mathematik
\usepackage{amsmath}			% Formeln, align-Umgebung usw.
\usepackage{amssymb}			% zusätzliche Symbole, z. B. \mathbb{R}
\usepackage{amsthm}				% Satz, Lemma, Beweis ...

%% -- Bilder
\usepackage{graphicx}			% \includegraphics

%% -- Literaturverzeichnis mit biblatex (Programm: biber)
\usepackage[style=numeric, backend=biber]{biblatex}
\addbibresource{\jobname.bib}	% die .bib-Datei (hier: oben in TEIL 0 erzeugt)

%% -- Anklickbare Links und Querverweise im PDF (als letztes Paket laden)
\usepackage[hidelinks]{hyperref}

%% -- Theorem-Umgebungen: \newtheorem{Name}[mitzählen]{Überschrift}[Nummerierung]
\theoremstyle{plain}				% kursiver Text
\newtheorem{satz}{Satz}[section]	% Satz 1.1, Satz 1.2, ...
\newtheorem{lemma}[satz]{Lemma}		% Lemma zählt mit den Sätzen mit
\theoremstyle{definition}			% aufrechter Text
\newtheorem{definition}[satz]{Definition}
\newtheorem{beispiel}[satz]{Beispiel}

%% -- Eigene Abkürzungen (Makros): sparen Tipparbeit
\newcommand{\R}{\mathbb{R}}		% \R ergibt die reellen Zahlen
\newcommand{\N}{\mathbb{N}}		% \N ergibt die natürlichen Zahlen

%% -- Titel, Autor, Datum
\title{Meine erste Ausarbeitung}
\subtitle{Proseminar Analysis}
\author{Vorname Nachname}
\date{\today}					% heutiges Datum; oder z. B. {24.~September 2026}

%% ====================================================================
%% TEIL 2: DER TEXT
%% Alles zwischen \begin{document} und \end{document} erscheint im PDF.
%% ====================================================================
\begin{document}

\maketitle			% setzt Titel, Autor, Datum
\tableofcontents	% Inhaltsverzeichnis (erst nach dem 2. Übersetzen vollständig)

%% --------------------------------------------------------------------
\section{Text schreiben}\label{sec:text}
%% --------------------------------------------------------------------

%% Tipp: Jeden Satz in eine eigene Zeile schreiben.
%% Ein einzelner Zeilenumbruch im Quelltext ist für LaTeX nur ein Leerzeichen.
Das ist der erste Satz.
Das ist der zweite Satz; er steht im PDF im selben Absatz.

%% Eine LEERZEILE beginnt einen neuen Absatz. Bitte NICHT \\ verwenden!
Das ist ein neuer Absatz.
Man kann Wörter \emph{hervorheben} und \enquote{Anführungszeichen} richtig setzen.
Fußnoten sind ganz einfach.\footnote{So sieht eine Fußnote aus.}

Striche gibt es in drei Längen:
den Bindestrich (Riemann-Integral, eingegeben als \verb|-|),
den Gedankenstrich -- wie hier -- und den Bis-Strich (Seite 10--20), beide eingegeben als \verb|--|,
sowie das Minuszeichen $-1$ im Mathemodus.

Abkürzungen bekommen einen kleinen Abstand: z.\,B., d.\,h., u.\,a.

\subsection{Aufzählungen}

Eine Liste mit Punkten:
\begin{itemize}
	\item erster Punkt
	\item zweiter Punkt
\end{itemize}

Eine nummerierte Liste:
\begin{enumerate}
	\item erstens
	\item zweitens
\end{enumerate}

%% --------------------------------------------------------------------
\section{Mathematik}\label{sec:mathe}
%% --------------------------------------------------------------------

\subsection{Formeln im Text und abgesetzt}

%% Formeln im Text stehen zwischen $ ... $
Eine Formel im Text: Für alle $x \in \R$ gilt $x^2 \geq 0$.
Brüche, Wurzeln und Indizes: $\frac{a}{b}$, $\sqrt{2}$, $a_n$, $x^{2n}$.

%% Abgesetzte Formel OHNE Nummer: \[ ... \]   (niemals $$ ... $$ verwenden!)
Eine abgesetzte Formel ohne Nummer:
\[
	\sum_{k=1}^{n} k = \frac{n(n+1)}{2} .
\]

%% Abgesetzte Formel MIT Nummer: equation-Umgebung mit \label
Eine abgesetzte Formel mit Nummer:
\begin{equation}\label{eq:euler}
	\mathrm{e}^{\mathrm{i}\pi} + 1 = 0 .
\end{equation} 
Auf sie verweist man mit \verb|\eqref{eq:euler}|, das ergibt \eqref{eq:euler}.

%% Mehrzeilige Formeln: align*-Umgebung; & markiert die Stelle, an der
%% ausgerichtet wird, \\ beendet eine Formelzeile (nur hier erlaubt!)
Mehrzeilige Rechnungen richtet man am Gleichheitszeichen aus:
\begin{align*}
	(a+b)^2 &= (a+b)(a+b) \\
	        &= a^2 + 2ab + b^2 .
\end{align*}

%% Große Klammern passen sich mit \left( ... \right) an den Inhalt an
Klammern um Brüche: $\left( \frac{1}{2} \right)^n$ statt $(\frac{1}{2})^n$.

\subsection{Sätze und Beweise}

\begin{definition}\label{def:beschraenkt}
Eine Folge $(a_n)_{n\in\N}$ heißt \emph{beschränkt}, wenn es ein $M > 0$ gibt mit $|a_n| \leq M$ für alle $n \in \N$.
\end{definition}

\begin{satz}\label{satz:konvergent}
Jede konvergente Folge ist beschränkt.
\end{satz}

\begin{proof}
Sei $(a_n)$ konvergent mit Grenzwert $a$.
Dann gibt es ein $N$ mit $|a_n - a| < 1$ für alle $n \geq N$, also $|a_n| < |a| + 1$.
Mit $M := \max\{|a_1|, \dots, |a_{N}|, |a| + 1\}$ folgt die Behauptung.
\end{proof}

%% Querverweis: ~ ist ein geschütztes Leerzeichen (dort wird nie umgebrochen)
Nach Satz~\ref{satz:konvergent} ist jede konvergente Folge im Sinne von Definition~\ref{def:beschraenkt} beschränkt.
Die Mathematik steht in Abschnitt~\ref{sec:mathe}.

%% --------------------------------------------------------------------
\section{Bilder und Tabellen}
%% --------------------------------------------------------------------

%% Bilder und Tabellen „gleiten“ – LaTeX setzt sie an eine günstige Stelle.
%% Immer mit \caption (Beschriftung) und DANACH \label (für Verweise).
Abbildung~\ref{fig:beispiel} zeigt ein Beispielbild, Tabelle~\ref{tab:werte} einige Werte.

\begin{figure}[htb]
	\centering
	%% eigenes Bild: Datei in denselben Ordner legen, Name ohne Endung angeben
	\includegraphics[width=0.4\textwidth]{example-image-a}
	\caption{Ein Beispielbild}
	\label{fig:beispiel}
\end{figure}

\begin{table}[htb]
	\centering
	\caption{Einige Quadratzahlen}	% bei Tabellen steht die Beschriftung oben
	\label{tab:werte}
	\begin{tabular}{c|c}			% zwei Spalten, zentriert, mit Trennstrich
		$n$ & $n^2$ \\ \hline
		1   & 1     \\
		2   & 4     \\
		3   & 9
	\end{tabular}
\end{table}

%% --------------------------------------------------------------------
\section{Literatur zitieren}
%% --------------------------------------------------------------------

%% Die Einträge stehen in der .bib-Datei (siehe TEIL 0 ganz oben).
%% Zitiert wird mit dem Schlüssel des Eintrags.
Ein gutes Buch zur Analysis ist \cite{heuser}.
Mit Seitenangabe: \cite[S.~42]{heuser}.
Zu LaTeX selbst empfiehlt sich \cite{lshort}.

%% Das Literaturverzeichnis wird automatisch erstellt (mit biber!)
\printbibliography

\end{document}
